\documentclass[11pt]{article}

\usepackage[margin=1in]{geometry}
\usepackage{amsmath,amssymb,amsthm}
\usepackage{hyperref}

\newtheorem{theorem}{Theorem}
\newtheorem{lemma}{Lemma}
\newtheorem{proposition}{Proposition}
\newtheorem{corollary}{Corollary}
\newtheorem{assumption}{Assumption}
\newtheorem{definition}{Definition}
\newtheorem{remark}{Remark}

\title{Hedge on a $1/\sqrt{T}$ Grid with a Concentrated Prior: Regret and Concentration}
\date{}
\author{}

\begin{document}
\maketitle

\section{Problem setup}

Fix a horizon $T\in\mathbb{N}$. Consider online learning over $[-1,1]$ using a uniform grid of spacing
\[
    h := \frac{2}{N}, \qquad N\in\mathbb{N}.
\]
The grid points are
\[
    x_i := -1 + i h,\qquad i=0,1,\dots,N.
\]
You may think of $h \asymp 1/\sqrt{T}$, i.e.\ $N\asymp \sqrt{T}$.

At each round $t=1,2,\dots,T$:
\begin{itemize}
    \item the learner chooses a distribution $a_t \in \Delta^{N+1}$ over the grid,
    \item samples an action $X_t\in\{x_0,\dots,x_N\}$ with $\mathbb{P}(X_t=x_i)=a_t(i)$,
    \item then observes a convex function $f_t:[-1,1]\to \mathbb{R}$ and incurs loss $f_t(X_t)$.
\end{itemize}

We will analyze Hedge on the finite action set $\{x_i\}$ using the losses evaluated on the grid.

\subsection*{Standing boundedness assumption (for regret)}
The standard Hedge regret bound is cleanest when losses are bounded. We assume:
\begin{assumption}[Bounded losses on the grid]\label{ass:bounded}
    For all $t$ and all grid points $x_i$, the grid losses satisfy
    \[
        0 \le \ell_t(i) := f_t(x_i) \le 1.
    \]
\end{assumption}
This can be enforced by scaling/clipping if needed, since only regret statements depend on it.

\subsection*{Convexity assumption (for log-concavity / concentration propagation)}
\begin{assumption}[Convex losses]\label{ass:convex}
    Each $f_t$ is convex on $[-1,1]$.
\end{assumption}

\section{Algorithm: Hedge with an arbitrary prior}

Fix a learning rate $\eta>0$ and an initial distribution (prior) $a_1\in\Delta^{N+1}$ with $a_1(i)>0$ for all $i$.

Hedge updates are
\begin{equation}\label{eq:hedge}
    a_{t+1}(i) \;=\; \frac{a_t(i)\exp\!\big(-\eta\,\ell_t(i)\big)}{\sum_{j=0}^N a_t(j)\exp\!\big(-\eta\,\ell_t(j)\big)}\,,\qquad i=0,\dots,N.
\end{equation}
Unrolling \eqref{eq:hedge} gives the closed form
\begin{equation}\label{eq:closedform}
    a_t(i)\ \propto\ a_1(i)\exp\!\Big(-\eta\sum_{s=1}^{t-1}\ell_s(i)\Big).
\end{equation}

\section{Regret with a non-uniform (concentrated) prior}

Define the expected (pseudo-)regret against the best \emph{grid} action:
\[
    \mathrm{Reg}^{\mathrm{grid}}_T
    := \sum_{t=1}^T \mathbb{E}\big[\ell_t(X_t)\big] - \min_{i\in\{0,\dots,N\}} \sum_{t=1}^T \ell_t(i),
\]
where the expectation is over the learner's randomization.

\begin{theorem}[Standard Hedge bound with an arbitrary prior]\label{thm:hedge-regret}
    Under Assumption~\ref{ass:bounded}, for any comparator distribution $u\in\Delta^{N+1}$,
    \[
        \sum_{t=1}^T \langle a_t,\ell_t\rangle - \sum_{t=1}^T \langle u,\ell_t\rangle
        \ \le\
        \frac{\mathrm{KL}(u\|a_1)}{\eta} + \frac{\eta T}{8},
    \]
    where $\langle a,\ell\rangle=\sum_i a(i)\ell(i)$ and $\mathrm{KL}(u\|a_1)=\sum_i u(i)\log\frac{u(i)}{a_1(i)}$.
    In particular, for any fixed action $i^\star$,
    \[
        \sum_{t=1}^T \mathbb{E}\big[\ell_t(X_t)\big] - \sum_{t=1}^T \ell_t(i^\star)
        \ \le\
        \frac{\log\!\big(1/a_1(i^\star)\big)}{\eta} + \frac{\eta T}{8}.
    \]
\end{theorem}

\begin{corollary}[Worst-case grid regret]\label{cor:worst-grid}
    Let $a_{\min}:=\min_i a_1(i)$. Then
    \[
        \mathrm{Reg}^{\mathrm{grid}}_T \ \le\ \frac{\log(1/a_{\min})}{\eta}+\frac{\eta T}{8}.
    \]
    Optimizing over $\eta$ yields
    \[
        \mathrm{Reg}^{\mathrm{grid}}_T \ \le\ \sqrt{\frac{T}{2}\,\log\!\frac{1}{a_{\min}}}\,.
    \]
\end{corollary}

\subsection{A concentrated prior that still preserves $\widetilde{O}(\sqrt{T})$ regret}

A convenient centered, concentrated initialization is the \emph{discrete Gaussian prior}:
\begin{equation}\label{eq:gaussian-prior}
    a_1(i) := \frac{\exp\!\big(-x_i^2/(2\sigma^2)\big)}{\sum_{j=0}^N \exp\!\big(-x_j^2/(2\sigma^2)\big)}\,,\qquad \sigma>0.
\end{equation}

The smallest prior mass occurs at the endpoints ($|x_i|=1$), hence
\[
    a_{\min} = \frac{\exp\!\big(-1/(2\sigma^2)\big)}{Z},\qquad
    Z:=\sum_{j=0}^N \exp\!\big(-x_j^2/(2\sigma^2)\big).
\]
Since each term in $Z$ is at most $1$, we have $Z\le N+1$ and therefore
\begin{equation}\label{eq:amin-bound}
    \log\frac{1}{a_{\min}}
    = \log Z + \frac{1}{2\sigma^2}
    \ \le\ \log(N+1)+\frac{1}{2\sigma^2}.
\end{equation}

If you choose the grid resolution $h\asymp 1/\sqrt{T}$, then $N\asymp \sqrt{T}$ and $\log(N+1)=\Theta(\log T)$.
Thus, choosing
\begin{equation}\label{eq:sigma-choice}
    \frac{1}{\sigma^2} = \Theta(\log T)
    \quad\Longleftrightarrow\quad
    \sigma^2 = \Theta\!\Big(\frac{1}{\log T}\Big)
\end{equation}
keeps $\log(1/a_{\min})=\Theta(\log T)$ and therefore (by Corollary~\ref{cor:worst-grid})
\[
    \mathrm{Reg}^{\mathrm{grid}}_T = O\!\big(\sqrt{T\log T}\big),
\]
which is $\widetilde{O}(\sqrt{T})$ (i.e.\ $\sqrt{T}$ up to logarithmic factors).

\begin{remark}
    Your prompt asked about keeping ``$\widetilde{O}(T)$'' regret, but the standard target in this setting is $\widetilde{O}(\sqrt{T})$ regret; the bound above achieves $\widetilde{O}(\sqrt{T})$ even with a non-uniform concentrated prior.
\end{remark}

\subsection{Regret against the best \emph{continuous} point in $[-1,1]$}

To compare against the best $x\in[-1,1]$, one needs a discretization condition; a standard one is Lipschitzness.

\begin{assumption}[Lipschitz losses (optional)]\label{ass:lipschitz}
    Each $f_t$ is $L$-Lipschitz on $[-1,1]$:
    \[
        |f_t(x)-f_t(y)|\le L|x-y| \quad \forall x,y\in[-1,1].
    \]
\end{assumption}

Let $x^\star\in[-1,1]$ minimize $\sum_{t=1}^T f_t(x)$. Let $\hat x^\star$ be its nearest grid point, so $|\hat x^\star-x^\star|\le h/2$.
Then by Lipschitzness,
\[
    \sum_{t=1}^T f_t(\hat x^\star) - \sum_{t=1}^T f_t(x^\star)\ \le\ \frac{LTh}{2}.
\]
Hence the regret against the best continuous point is
\[
    \mathrm{Reg}^{\mathrm{cont}}_T
    := \sum_{t=1}^T \mathbb{E}[f_t(X_t)] - \min_{x\in[-1,1]}\sum_{t=1}^T f_t(x)
    \ \le\
    \mathrm{Reg}^{\mathrm{grid}}_T + \frac{LTh}{2}.
\]
With $h=\Theta(1/\sqrt{T})$, this discretization term is $O(L\sqrt{T})$, matching the $\sqrt{T}$ scale.

\section{Log-concavity and ``strong log-concavity'' on a grid}

Define the (discrete) potential of a distribution $a\in\Delta^{N+1}$ by
\[
    V(i) := -\log a(i).
\]
For any sequence $g(i)$ on $\{0,\dots,N\}$ define the interior second difference
\[
    \Delta^2 g(i) := g(i+1)-2g(i)+g(i-1),\qquad i=1,\dots,N-1.
\]
A sequence $g$ is \emph{discrete convex} if $\Delta^2 g(i)\ge 0$ for all interior $i$.

\begin{definition}[Grid log-concavity and grid strong log-concavity]\label{def:slc}
    A distribution $a$ on the grid is:
    \begin{itemize}
        \item \emph{log-concave on the grid} if $i\mapsto \log a(i)$ is discrete concave, equivalently $i\mapsto V(i)$ is discrete convex: $\Delta^2 V(i)\ge 0$.
        \item \emph{$m$-strongly log-concave on the grid} (for $m\ge 0$) if
              \[
                  \Delta^2 V(i) \ \ge\ m h^2\qquad \text{for all } i=1,\dots,N-1.
              \]
    \end{itemize}
\end{definition}

The scaling by $h^2$ is chosen so that $m$ is stable as $h\to 0$; for smooth $V$, one has $\Delta^2 V(i)\approx h^2 V''(x_i)$.

\subsection{Convex losses preserve (strong) log-concavity}
Let $V_t(i):=-\log a_t(i)$.

\begin{proposition}[Curvature monotonicity under convex losses]\label{prop:kappa-mono}
    Assume $f_t$ is convex (Assumption~\ref{ass:convex}). Then the grid sequence $i\mapsto \ell_t(i)=f_t(x_i)$ is discrete convex, i.e.\ $\Delta^2 \ell_t(i)\ge 0$.
    Moreover, the Hedge update satisfies
    \[
        V_{t+1}(i) = V_t(i) + \eta\,\ell_t(i) + \text{(constant independent of $i$)}.
    \]
    Consequently, if $a_t$ is $m$-strongly log-concave on the grid, then so is $a_{t+1}$ (with the same $m$). In particular, $m$ is non-decreasing over time.
\end{proposition}

\begin{proof}
    From the update \eqref{eq:hedge},
    \[
        \log a_{t+1}(i)=\log a_t(i)-\eta\,\ell_t(i)-\log\Big(\sum_j a_t(j)e^{-\eta\ell_t(j)}\Big),
    \]
    so negating gives
    \[
        V_{t+1}(i)=V_t(i)+\eta\,\ell_t(i)+\log Z_t
    \]
    where $\log Z_t$ is constant in $i$. Taking second differences kills the constant:
    \[
        \Delta^2 V_{t+1}(i)=\Delta^2 V_t(i)+\eta\,\Delta^2 \ell_t(i).
    \]
    If $f_t$ is convex on $[-1,1]$, then its samples on an equally spaced grid form a discrete convex sequence, hence $\Delta^2 \ell_t(i)\ge 0$.
    Thus $\Delta^2 V_{t+1}(i)\ge \Delta^2 V_t(i)$ pointwise, proving monotonicity and preservation of any lower bound $\Delta^2 V_t(i)\ge m h^2$.
\end{proof}

\subsection{Discrete Gaussian prior is $m$-strongly log-concave with $m=1/\sigma^2$}

\begin{proposition}[Curvature of the discrete Gaussian prior]\label{prop:gauss-curv}
    Let $a_1$ be the discrete Gaussian prior \eqref{eq:gaussian-prior}. Then $a_1$ is $m$-strongly log-concave on the grid with
    \[
        m=\frac{1}{\sigma^2}.
    \]
    Consequently, under convex losses, every iterate $a_t$ is $m$-strongly log-concave with the same $m=1/\sigma^2$.
\end{proposition}

\begin{proof}
    For the prior \eqref{eq:gaussian-prior},
    \[
        V_1(i)=-\log a_1(i)=\frac{x_i^2}{2\sigma^2}+\text{const}.
    \]
    Compute the second difference using $x_{i\pm 1}=x_i\pm h$:
    \[
        \Delta^2 V_1(i)
        = \frac{(x_i+h)^2 - 2x_i^2 + (x_i-h)^2}{2\sigma^2}
        = \frac{2h^2}{2\sigma^2}
        = \frac{h^2}{\sigma^2}.
    \]
    Thus $\Delta^2 V_1(i)\equiv (1/\sigma^2)h^2$ for all interior $i$, i.e.\ $m=1/\sigma^2$.
    The final claim follows from Proposition~\ref{prop:kappa-mono}.
\end{proof}

\section{Concentration around the mean at every round}

Fix $t$ and let $X_t$ be the grid-valued random variable with $\mathbb{P}(X_t=x_i)=a_t(i)$.
Let
\[
    \mu_t := \mathbb{E}[X_t].
\]

There are (at least) two standard ways to convert strong log-concavity into mean-centered concentration:
\begin{itemize}
    \item A Poincar\'e/variance bound $\mathrm{Var}(X_t)\lesssim 1/m$, giving Chebyshev-type tails.
    \item A log-Sobolev/moment-generating-function bound giving sub-Gaussian (Chernoff) tails with variance proxy $\asymp 1/m$.
\end{itemize}

Because you requested a clean $(\epsilon,\delta)$-style statement, we record the following widely used implication of strong log-concavity (on a line) in a form suited to this note.

\begin{assumption}[Mean-centered sub-Gaussianity from $m$-strong log-concavity]\label{ass:subg}
    There exists an absolute constant $c_0>0$ such that if a grid distribution is $m$-strongly log-concave (Definition~\ref{def:slc}), then for $X$ drawn from that distribution,
    \[
        \mathbb{E}\exp\!\big(\lambda(X-\mathbb{E}X)\big)\ \le\ \exp\!\Big(\frac{\lambda^2}{2c_0 m}\Big)
        \qquad\forall \lambda\in\mathbb{R}.
    \]
    Equivalently (by Chernoff bounds), for all $\epsilon>0$,
    \[
        \mathbb{P}\big(|X-\mathbb{E}X|\ge \epsilon\big)\ \le\ 2\exp\!\big(-c_0 m \epsilon^2\big).
    \]
\end{assumption}

\begin{remark}
    Assumption~\ref{ass:subg} is the discrete-grid analogue of the standard continuous fact: if a density is proportional to $\exp(-V(x))$ with $V''(x)\ge m$, then it satisfies a log-Sobolev inequality with constant $1/m$ and hence is sub-Gaussian with variance proxy $1/m$. On a 1D grid (a birth--death geometry), analogous (modified) log-Sobolev and concentration statements are available for strongly log-concave measures. This note treats the sub-Gaussian tail bound as a black box so that the dependence on $(T,h,\sigma)$ is fully explicit.
\end{remark}

Under this assumption, concentration for the Hedge iterates is immediate.

\begin{theorem}[Uniform-in-time concentration for Hedge with a discrete Gaussian prior]\label{thm:concentration}
    Assume convex losses (Assumption~\ref{ass:convex}) and initialize Hedge with the discrete Gaussian prior \eqref{eq:gaussian-prior}.
    Then for every round $t$, the policy $a_t$ is $m$-strongly log-concave with $m=1/\sigma^2$, and therefore (by Assumption~\ref{ass:subg}) satisfies the mean-centered tail bound
    \begin{equation}\label{eq:mean-tail}
        \mathbb{P}\big(|X_t-\mu_t|\ge \epsilon\big)\ \le\ 2\exp\!\Big(-c_0\,\frac{\epsilon^2}{\sigma^2}\Big)
        \qquad\forall \epsilon>0,
    \end{equation}
    uniformly over all $t=1,\dots,T$.
    Equivalently, for any $\delta\in(0,1)$,
    \begin{equation}\label{eq:eps-delta}
        \mathbb{P}\Big(|X_t-\mu_t|\le \sigma\sqrt{\frac{1}{c_0}\log\frac{2}{\delta}}\Big)\ \ge\ 1-\delta
        \qquad\forall t.
    \end{equation}
\end{theorem}

\begin{proof}
    By Proposition~\ref{prop:gauss-curv} and Proposition~\ref{prop:kappa-mono}, $a_t$ is $m$-strongly log-concave with $m=1/\sigma^2$ for all $t$. Apply Assumption~\ref{ass:subg}.
\end{proof}

\subsection{How many grid ``ticks'' contain most of the mass?}

Fix $\delta\in(0,1)$. Define the high-probability radius
\[
    \epsilon(\delta) := \sigma\sqrt{\frac{1}{c_0}\log\frac{2}{\delta}}.
\]
Then \eqref{eq:eps-delta} implies that, for every $t$, at least $(1-\delta)$ of the probability mass lies in the interval
\[
    [\mu_t-\epsilon(\delta),\ \mu_t+\epsilon(\delta)].
\]
The number of grid points in an interval of length $2\epsilon$ is at most $2\lceil \epsilon/h\rceil+1$.
Hence, for every $t$,
\begin{equation}\label{eq:ticks}
    \mathbb{P}\Big(X_t \in \{x_i:\ |x_i-\mu_t|\le \epsilon(\delta)\}\Big)\ \ge\ 1-\delta,
\end{equation}
and the set $\{x_i:\ |x_i-\mu_t|\le \epsilon(\delta)\}$ contains at most
\begin{equation}\label{eq:tickcount}
    M(\delta)\ :=\ 2\left\lceil \frac{\epsilon(\delta)}{h}\right\rceil + 1
    \ =\ 2\left\lceil \frac{\sigma}{h}\sqrt{\frac{1}{c_0}\log\frac{2}{\delta}}\right\rceil + 1
\end{equation}
grid points.

If you set $h=\Theta(1/\sqrt{T})$ and $\sigma^2=\Theta(1/\log T)$ as in \eqref{eq:sigma-choice}, then
\[
    \frac{\sigma}{h}\ =\ \Theta\!\Big(\sqrt{\frac{T}{\log T}}\Big),
\]
so for any fixed confidence level (e.g.\ constant $\delta$), the number of ticks capturing at least $(1-\delta)$ mass scales as
\[
    M(\delta)\ =\ \Theta\!\Big(\sqrt{\frac{T}{\log T}}\Big)\cdot \sqrt{\log\frac{1}{\delta}}
    \quad\text{(up to absolute constant factors and the ceiling)}.
\]
In particular, it is \emph{not} a constant number of ticks when $h\asymp 1/\sqrt{T}$.

\section{Putting it together: one concrete choice of parameters}

Choose the grid resolution
\[
    h := \Theta\!\Big(\frac{1}{\sqrt{T}}\Big)\qquad (N=\Theta(\sqrt{T})).
\]
Initialize Hedge with the discrete Gaussian prior \eqref{eq:gaussian-prior} with
\[
    \sigma^2 := \frac{1}{\alpha \log T}\qquad \text{for some fixed }\alpha>0.
\]
Then:
\begin{itemize}
    \item \textbf{Regret.} By \eqref{eq:amin-bound}, $\log(1/a_{\min}) \le \log(N+1)+\alpha \log T/2 = \Theta(\log T)$.
          Thus Hedge attains $\mathrm{Reg}^{\mathrm{grid}}_T = O(\sqrt{T\log T})=\widetilde{O}(\sqrt{T})$.
          If in addition Assumption~\ref{ass:lipschitz} holds, the regret against the best continuous point is
          \[
              \mathrm{Reg}^{\mathrm{cont}}_T \ \le\ O(\sqrt{T\log T}) + O(L\sqrt{T}) \ =\ \widetilde{O}(\sqrt{T}).
          \]
    \item \textbf{Concentration at every round.} For every $t$, $X_t$ is concentrated around its mean with
          \[
              \mathbb{P}\big(|X_t-\mu_t|\ge \epsilon\big)\ \le\ 2\exp\!\big(-c_0 \alpha (\log T)\,\epsilon^2\big).
          \]
          Equivalently, with probability at least $1-\delta$,
          \[
              |X_t-\mu_t|\ \le\ \sqrt{\frac{1}{c_0\alpha}\cdot \frac{\log(2/\delta)}{\log T}}.
          \]
          The corresponding number of grid ticks capturing at least $(1-\delta)$ mass is at most
          \[
              M(\delta)
              \ =\
              2\left\lceil
              \sqrt{\frac{T}{c_0\alpha\log T}}\ \sqrt{\log\frac{2}{\delta}}
              \right\rceil + 1.
          \]
\end{itemize}

\section{Tradeoff note (why you cannot generally get ``constant ticks'' and worst-case $\widetilde{O}(\sqrt{T})$ regret)}

If you wanted most mass to lie within $O(1)$ ticks, you would need $\epsilon(\delta)=O(h)$, i.e.\ $\sigma=O(h)$.
With $h\asymp 1/\sqrt{T}$, that means $\sigma^2=O(1/T)$ and hence $1/(2\sigma^2)=\Omega(T)$.
But then \eqref{eq:amin-bound} gives $\log(1/a_{\min})=\Omega(T)$, which makes the Hedge regret bound (Corollary~\ref{cor:worst-grid}) linear in $T$.
Thus, in the worst-case (adversarial) setting over all comparators in $[-1,1]$, there is an inherent tension:
very strong per-round concentration (at the grid scale) forces a large prior penalty and destroys $\widetilde{O}(\sqrt{T})$ regret guarantees.

\end{document}
